\chapter{About This Notebook}

\section{The Use of \LaTeX{}} This notebook was created using \LaTeX{}, a free,
open source typesetting program that is especially useful for technical
documentation \cite{bib:latex}.  We chose to use \LaTeX{} for the following
reasons:
\begin{itemize}
      \item \LaTeX{} provides greater control over the style and format of
            documents than traditional word processors. This is especially true
            for:
            \begin{itemize}
                  \item Long documents
                  \item Technical documents
            \end{itemize}
      \item Because \LaTeX{} documents are expressed with plain text, it
            provides an avenue for future automated documentation tools.
      \item Most practically, \LaTeX{} is a useful tool career-wise, so this
            notebook is a good opportunity to get practice with the software.
\end{itemize}
This notebooks format is inspired heavily by a template provided by Purdue
SIGBots \cite{bib:purduenotebook} and an open source controls engineering
textbook by Tyler Veness \cite{bib:controlinfrc} (which will probably be
referenced again).
\section{Formatting}
This notebook is presently divided into several parts which will now be
described in more detail.
\begin{itemize}
      \item \textit{Preface} deals with the notebook, team and club procedures,
            logistics, overall structure of our program, general game analysis,
            and design philosophy.
      \item \textit{Hardware} deals with the design process for the robots.
      \item \textit{Software} deals with the design process for software used
            with both robots as well as potentially custom electronics.
      \item \textit{Appendices} contain information that would not fit nicely in
            other sections of the notebook, such as commit logs, meeting notes,
            etcetera.
\end{itemize}
The sections regarding the hardware and software contain largely  "highlights"
of the team/club meetings. For a more traditional notebook experience, the
meeting note section of the appendix and the "logs" in each chapter should be
especially considered. This notebook's version history was recorded using git
and was distributed using GitHub. \textbf{To ensure that the notebook was made
      alongside the rest of the project, the commit history of the notebook repository
      should be attached. Additionally, the following QR code leads to the notebook
      repository.}

\pagebreak
\begin{center}
      \
      \vfill
      \qrcode[height = .6\linewidth]{https://github.com/The-ATU-Robotics-Club/ATUM-V5RC-High-Stakes-Notebook}

      \vspace*{40pt}
      \url{https://github.com/The-ATU-Robotics-Club/ATUM-V5RC-High-Stakes-Notebook}

      \textit{Leads to the GitHub repository for this notebook, it should be public for this competition.}

      \textit{If we have forgotten to change the visibility of the repository, the commit history should also be attached.}

      The PDF version may be preferred for reading, as the table of contents
      provides hyperlinks to various sections you will have to download the
      file, though. See the "releases" for the current PDF file.
      \vfill
\end{center}
\pagebreak

\addtocontents{toc}{\setcounter{tocdepth}{1}}
\section{Log}
This section deals with major changes to the format or workflow of this notebook
(not new entries). If the format of this notebook changes substantially, or a
new tool is implemented in its development, it should be mentioned here. These
changes are quite different from the rest of the design process, so notes of
this variety are isolated here.
\subsection{6/13/2023}
Set up beginning of notebook and the backbone for future entries. Add a command
for inserting figures, \sout{such as:} (removed; see 8/9/2024)
\subsection{6/14/2023}
Finalize basic setup and style for the notebook.
\subsection{6/15/2023}
Began using LTeX \cite{bib:ltex} extension for spell and grammar checks. Reset
chapter number with each part.
\subsection{6/16/2023}
Created a python script for generating \LaTeX{} for design matrices and a
corresponding command for further formatting, \sout{such as:} (removed; see 8/9/2024)
\subsection{6/19/2023}
Created a GitHub action for adding club meeting notes from the common
organization repository.

While this notebook otherwise records all meetings and decisions of the team,
adding club notes may provide greater context to this project history.
Additionally, this undertaking was a good learning experience for generating
\LaTeX{} files for further automation as well as gaining access to other private
organization repositories in GitHub actions. The initial solution for the latter
involved using a Personal Authentication Token (PAT) which is akin to sharing
your password around, insecure. The final solution involves using Deploys Keys
and SSH which is much more secure.
\subsection{6/29/2023}
Add an automatic workflow for adding changes to the GitHub Project to the
appendix. A primary goal for this notebook is to centralize organization and
documentation, so that a comprehensive record of the project is provided. This
workflow will ensure we do not forget to update the status of our project board
and timeline in the notebook.
\subsection{7/6/2023}
Add an easy-to-use workflow for adding version history information to the
notebook. Not only is this necessary to prove that the notebook was written
alongside the project, it also helps to provide a more complete picture of
project history.
\subsection{9/25/2023}
Improve formatting of chapters and section numbers. Prevent floating objects
from moving too far away from preferred location.
\subsection{10/29/2023}
Add several resources for how to add content to the notebook. Also add a GitHub
Workflow to the organization repository to automatically close issues under the
"Done" category of the project page.
\subsection{1/30/2024}
The notebook now has specific formatting for code snippets:
\begin{lstlisting}[language=Python]
    print("Hello world! Here's a number: " + 123); # This is a comment.
\end{lstlisting}
\subsection{2/5/2024}
Reformat how images appear in the notebook. Add special a format for inserting
pdfs. Combine the 15" and 24" parts into one hardware part. Greatly decrease
length of appendix by reducing font size. Improve the table of contents layout.
\subsection{7/28/2024}
We have decided to switch from using VSCode (with extensions) and GitHub for
writing the \LaTeX{} for the notebook, to using an online editor called Overleaf
\cite{bib:overleaf}. This was a pretty clear cut decision, with benefits
including:
\begin{itemize}
      \item The ability to work concurrently on the notebook with no risk of
            conflicts.
      \item Much easier set up with no download (only need to make an account).
      \item Macros and buttons for common tasks like creating sections/chapters,
            bolding/italicizing text, making tables, inserting images, and more.
      \item Integration with GitHub so all automation features still work.
      \item A visual editor to get immediate feedback on changes.
      \item Built in spellcheck (and easier dictionary support) and text wrapping.
      \item And much more...
\end{itemize}
It really is just better in every way. It does require a subscription, but we
unanimously considered this a worthy expense. Work in VSCode will pretty much
be relegated to working with the GitHub workflows.
\subsection{7/31/2024}
Improved the look of the commit log and gave up on automating changing GitHub
usernames to names. It was too inconvenient to collect it from members, so it'll
probably be more convenient to find and replace them manually if we want to.
\par Also improved the look of the notebook by changing the spacing between
sections and paragraphs.
\subsection{8/3/2024}
Improved the look of the project summary, shrank its size, and made it automatically
color the rows corresponding to task status.
\subsection{8/6/2024}
Added a workflow for quickly adding all the files from our code repository into the
notebook. Bringing all the files over manually before was laborsome and wasted a lot
of time that could have been dedicated to elaborating on other topics.
\subsection{8/7/2024}
Greatly improve the formatting for the team profile. It used to be just a list of
names, photos, and biographies that went on far too long. Significantly shortened
and better looking.
\subsection{8/9/2024}
Commit to the "floating" nature of images; removing the previous command for inserting
images and working on the formatting for figures and cross references. Similar changes,
such as removing several commands like insertimage or insertmatrix, were removed in favor
of encouraging team members to employ the Overleaf provided macros and simplify the
learning curve for the notebook. Removed the material specific to the previous year.
\subsection{9/12/2024}
Rework the file structure of the notebook to make it much easier for people to get started
with the notebook. 
\addtocontents{toc}{\setcounter{tocdepth}{\defaulttocdepth}}

\chapter{Technology \& Methodology}
\section{git/GitHub}
\section{Discord}
\section{Club and Team Organization}

\chapter{Team Profile}
\leftsidemember{Braden Pierce}{images/members/braden.jpg}{Change the
    world...}{ I am a programmer for the team, working primarily on the lower
    end side of the software; creating the tools and systems for developing our
    controls and autonomous routines, rather for others on the team to work on
    those aspects of the code. Additionally, I am the club Secretary,
    responsible for helping organize the club; from setting up and helping
    direct meetings, creating and distributing merchandise, and helping develop
    the tools we use to communicate and get things done. I am a senior double
    majoring in Computer Science and Mathematics. I have been competing in VEX
    Robotics for about seven years and it has been an excellent motivator for me
    to learn more regarding computer science, mathematics, and engineering. }